\chapter{Appendix}

\section*{Modular arithmetic}

\begin{theorem}
    (Division).

    For all $a, n \in \Z$ with $a > n$ there exist $q \in \Z$ and $r \in \Z \cap [0, q)$ such that $\frac{a}{n} = q + \frac{r}{n}$. The number $q$ is called the \textit{quotient}, and the number $r$ is called the \textit{remainder}. The remainder involved in the division of $a$ by $n$ is also called \textit{$a$ modulo $n$}, and is denoted $a \bmod n$. Therefore $\frac{a}{n} = q + \frac{a \bmod n}{n}$.
\end{theorem}

\begin{proof}    
    ...reference some standard textbook...
\end{proof}

\begin{example}
    (Modulo cycles).

    \begin{align*}
        0 \bmod 4 = 0 \\
        1 \bmod 4 = 1 \\
        2 \bmod 4 = 2 \\
        3 \bmod 4 = 3 \\
        4 \bmod 4 = 0 \\
        5 \bmod 4 = 1 \\
        6 \bmod 4 = 2 \\
        7 \bmod 4 = 3 \\
        8 \bmod 4 = 0
    \end{align*}
\end{example}

\begin{remark}
    (Moduli of negative integers).
    
    Let $n \in \Z$ be a positive integer. Take care to remember that computing the modulus of a negative integer relative to $n$ may seem counterintuitive, since for $a, n \in \Z$ with $a > 0$ then usually\footnote{Specifically, this holds for all $a$ with $a \bmod n \neq 0$.} $(-a) \bmod n \neq -(a \bmod n)$.

    Let's think about $-a \bmod n$ for positive $a \in \Z$. From above, we have $\frac{a}{n} = q + \frac{a \bmod n}{n}$. That is, $a = nq + r = nq + a \bmod n$. From this expression, we can clearly see that $nq$ is the greatest integer multiple of $n$ that is less than $a$. Since $a \bmod n = nq - a$, we have

    \begin{align*}
        a \bmod n = (\text{greatest integer multiple of $n$ that is less than $a$}) - a.
    \end{align*}

    Even more explicitly,

    \begin{align*}
        a \bmod n = 
        \begin{cases}
            (\text{greatest integer multiple of $n$ that is less than $a$}) - |a| & a \geq 0 \\
            (\text{greatest integer multiple of $n$ that is less than $a$}) + |a| & a < 0
        \end{cases}.
    \end{align*}

    When $a > 0$, the greatest integer multiple of $n$ that is less than $a$ is closer to $0$ than $a$, and when $a < 0$, this same multiple is further from $0$ than $a$. So you might say that computing $a \bmod n$ involves undershooting when $a > 0$ and overshooting when $a < 0$.
\end{remark}

The division theorem, and therefore the modulo concept, can be extended to the case when the numbers being divided are both real numbers.

\begin{theorem}
    (Extended division).

    For all $x, T \in \R$ with $x > T$ there exist $q \in Z$ and $x \bmod T \in [0, q)$ such that $\frac{x}{T} = q + \frac{x \bmod T}{T}$. As before, $q$ is called the \textit{quotient}, and $x \bmod T$, which is read as \textit{$x$ modulo $T$}, is called the \textit{remainder}.
\end{theorem}

The following theorem making use of a modulus obtained from extended division comes in useful in Chapter \ref{ch::lin_alg} when investigating angle.

\begin{theorem}
\label{ch::appendix::thm::negation_of_angle}
    (Negation of an angle).

    If $t \in [0, 2\pi)$ then $-t \bmod 2\pi = 2\pi - t$.
\end{theorem}

\begin{proof}
    If $t \in [0, 2\pi)$ then $-t \in (-2\pi, 0]$ and so the greatest integer multiple of $2\pi$ less than $-t$ is $-2\pi$. The modulus $-t \bmod 2\pi$ is the difference between $-t$ and this multiple, i.e. $-t - (-2\pi) = 2\pi - t$.
\end{proof}

\section*{Some precursors to category theory}

\subsection*{Another perspective on dual spaces}

In Chapter \ref{ch::motivated_intro}, we motivated the definition of the dual space by seeing how it became useful when decomposing a linear function between vector spaces into a sum of ``elementary compositions''. Then, we kind of stumbled onto the theorems that describe in what sense a dual space is actually dual.

Here we present a derivation of dual spaces that emphasizes the core duality idea from the start.

\begin{deriv}
\label{ch::appendix::deriv::function_input_duality}
    (Function-input duality).
    
    Let $V$ and $W$ be vector spaces over a field $K$, let $\vv \in V$, let $\ff:V \rightarrow W$ be a function, and consider the expression
    
    \begin{align*}
    	\ff(\vv).
    \end{align*}

    Typically, we think of $\ff$ as being applied to $\vv$ to produce the output $\ff(\vv)$. We can also think of $\vv$ as being applied to $\ff$, in some sense, to produce $\ff(\vv)$. In intuitive notation, we have $\text{``}\vv(\ff)\text{''} = \ff(\vv)$.

    To formalize the intuition, notice that for every $\vv \in V$ there is a function $\FF_{\vv}$ that itself takes a function $V \rightarrow W$ as input and applies it to $\vv$; so, $\FF_{\vv}(\ff) = \ff(\vv)$. Since the function sending $\vv \mapsto \FF_{\vv}$ is an injection, we can identify each $\vv \in V$ with $\FF_{\vv}$ and then think of $\vv$ as being applied to $\ff$ via the action of $\FF_{\vv}$. The intuition thus formalizes as $\text{``}\vv(\ff)\text{''} = \FF_{\vv}(\ff) = \ff(\vv)$.

	\vspace{.25cm}

    There is still more insight to be had! Appraising the situation as a whole, we see that the two functions involved in the two interpretations of the expression $\ff(\vv)$ are
    
    \begin{align*}
    	&\ff:V \rightarrow W \text{, acting on $\vv$}\\
    	&\FF_{\vv}:\{\text{functions $V \rightarrow W$}\} \rightarrow W \text{, acting on $\ff$}.
   	\end{align*}
   	
   	To make the notation easier, we define $V^* := \{\text{functions $V \rightarrow W$}\}$. The above can then be restated as
   	
   	\begin{align*}
   		&\ff:V \rightarrow W \text{, acting on $\vv$} \\
   		&\FF_{\vv}:V^* \rightarrow W \text{, acting on $\ff$}.
   	\end{align*}

	By defining the notation $V^*$, we have actually stumbled onto more algebraic structure. You might have to think about it for a second, but the above can \textit{further} be restated as

	\begin{align*}
		&\ff \in V^* \text{, acting on $\vv$} \\
		&\FF_{\vv} \in (V^*)^* \text{, acting on $\ff$}
	\end{align*}
	 
	 Overall, the relationship between $V$, $V^*$, and $(V^*)^*$ is that elements of $V \cong (V^*)^*$ and $V^*$ act on each other. Since elements of $V^*$ cannot be naturally identified with elements of $V$, yet can still act on elements of $V$, we call $V^*$ the \textit{dual set (to $V$)}. (We call $(V^*)^*$ the \textit{double dual set (to $V$)}), and denote it by $V^{**}$ for brevity.
	 
	 Except... there's a bit more to be done. We know $V$ can be injected within $(V^*)^*$, but we don't yet know $(V^*)^*$ can be injected within $V$. And it doesn't feel right to say we have ``duality'' \textit{unless} we also have symmetry. 
	 
	 So, when \textit{do} we have a natural isomorphism $(V^*)^* \cong V$? Well, making use of the natural isomorphisms from Theorem \ref{ch::motivated_intro::thm::four_fundamental_isos}, we have
	
	\begin{align*}
		(V^*)^* &:= \LLLL(\LLLL(V \rightarrow W) \rightarrow W) \\ &\cong \LLLL(V \rightarrow W) \otimes W^* \\
		&\cong V \otimes W^* \otimes W^* \\
		&\cong V \otimes (W^*)^{\otimes 2}.
	\end{align*}
	
	Thus $(V^*)^* \cong V$ if and only if $(W^*)^{\otimes 2} \cong K$. And this is the case if and only if $W \cong K$. The simplest way to achieve this is by simply setting $W := K$. We therefore update our definition:
	
	\begin{align*}
		V^* := \LLLL(V \rightarrow K)
	\end{align*} 
	
	Now, not only can elements of $V$ be represented by elements of $V^{**}$, but also can elements of $V^{**}$ be represented by elements of $V$.
\end{deriv}

\subsection*{Characterizing natural linear isomorphisms}

Let $V$ be a vector space. We know that $V$ is linearly isomorphic to $V^*$, since, if we choose a basis $E = \{\ee_1, ..., \ee_n\}$ for $V$, then there is an induced dual basis $E^* = \{\phi^{\ee_1}, ..., \phi^{\ee_n}\}$ for $V^*$, and the function $V \rightarrow V^*$ defined by $\ee_i \mapsto \phi^{\ee_i}$ is a linear isomorphism.

We also intuitively know that since this isomorphism depends on the basis $E$, it is dependent on circumstance, and not truly \textit{natural}. We haven't yet formalized this notion of naturality, though.

We do this now.

\begin{deriv}
	(Characterizing natural linear isomorphisms between $V$ and $V^*$).
	
	Heuristically, we posit that a linear isomorphism $\ff_V:V \rightarrow V^*$ should only be considered \textit{natural} if ``its manner of construction is preserved'' when the basis of $V$ is changed, or when vectors in $V$ are otherwise swapped out for identical vectors from a different vector space $W$ with a linear isomorphism $\gg:V \rightarrow W$. To make this heuristic precise, we will say that the manner of construction of $\ff_V$ is preserved if and only if the isomorphism corresponding to $\ff_V$ between $W$ and $W^*$ \textit{that is naturally induced by $\gg$} is equal to $\ff_W$.
	
	We have shifted the question ``What does it mean for a linear isomorphism to be natural?'' to the question ``What is the isomorphism corresponding to $\ff_V$ between $W$ and $W^*$ that is naturally induced by $\gg$?''. To answer this later question, we need to figure out what functions should label the unlabeled arrows in the following diagram:
	
	% https://tikzcd.yichuanshen.de/#N4Igdg9gJgpgziAXAbVABwnAlgFyxMJZARgBoAGAXVJADcBDAGwFcYkQB1EAX1PU1z5CKchWp0mrdgDUefEBmx4CRUcXEMWbRCGkA9AFRz+SoUTLqamqTo6Ge4mFADm8IqABmAJwgBbJKIgOBBIxLyePv6IZEEhiABM4SDefkjxNMFIAMzclNxAA
	\begin{center}
		\begin{tikzcd}
			V \arrow[d, "\ff_V"']   & W \arrow[l] \\
			V^* \arrow[r] & W^*
		\end{tikzcd}
	\end{center}
	
	It is hopefully reasonable to take as axiom that the natural way to go between $V$ and $W$ is to use $\gg:V \rightarrow W$ (or its inverse), as $\gg$ and $W$ were introduced in the same breath. And, since each linear function between vector spaces corresponds to a dual linear function between dual vector spaces, it is hopefully reasonable to take as axiom that the natural way to go between the dual spaces $V^*$ and $W^*$ is to use the dual $\gg^*:W^* \rightarrow V^*$ (or its inverse). If we do accept these axioms, then the completed diagram is
	
	% https://tikzcd.yichuanshen.de/#N4Igdg9gJgpgziAXAbVABwnAlgFyxMJZARgBoAGAXVJADcBDAGwFcYkQB1EAX1PU1z5CKchWp0mrdgDUefEBmx4CRUcXEMWbRCGkA9AFRz+SoUTLqamqTo6Ge4mFADm8IqABmAJwgBbJKIgOBBIZBJa7AA6kc7OesAAtMTcIDSM9ABGMIwACgLKwiBeWM4AFjjGIN5+oTTBSABMVpLaINEeHgD6smmZ2XmmKjrFZRW8nj7+iE1BIYgAzM0ROgAU0bGGAJTxSSm9Wbn5ZsMl5Q7cQA
	\begin{center}
		\begin{tikzcd}
			V \arrow[d, "\ff_V"']          & W \arrow[l, "\gg^{-1}"'] \\
			V^* \arrow[r, "(\gg^*)^{-1}"'] & W^*
		\end{tikzcd}
	\end{center}
	
	Now recall that we said $\ff_V$ should only be considered a natural linear isomorphism if and only if the isomorphism corresponding to $\ff_V$ between $W$ and $W^*$ that is naturally induced by $\gg$ is equal to $\ff_W$. This means that $\ff_V$ is a natural linear isomorphism if and only if for all vector spaces $W$ and linear isomorphisms $\gg:V \rightarrow W$, the following diagram commutes:
	
	% https://tikzcd.yichuanshen.de/#N4Igdg9gJgpgziAXAbVABwnAlgFyxMJZABgBpiBdUkANwEMAbAVxiRADUQBfU9TXfIRQBGclVqMWbAOrdeIDNjwEiZYePrNWiDgD0AVHL5LBRUeuqapO6Qe7iYUAObwioAGYAnCAFskAZmocCCQAJktJbRAAHWj3dzseD28-RDIQYKRRCS02WKcnIxAvXyR0zMRwkAY6ACMYBgAFfmUhEE8sJwALHBAI3J1Y+IB9TiTilKygkMRAnOsY6ILR+y4gA
	\begin{center}
		\begin{tikzcd}
			V \arrow[r, "\gg"] \arrow[d, "\ff_V"'] & W \arrow[d, "\ff_W"]   \\
			V^*                                    & W^* \arrow[l, "\ff^*"]
		\end{tikzcd}
	\end{center}
\end{deriv}

To be clear, we restate the definition we just derived.

\begin{defn}
	(Natural linear isomorphism between $V$ and $V^*$).
	
	When $V$ is a vector space, we say that a linear isomorphism $\ff_V:V \rightarrow V^*$ is \textit{natural} if and only if the following diagram commutes for all vector spaces $W$ and all linear isomorphisms $\gg:V \rightarrow W$:
	
	% https://tikzcd.yichuanshen.de/#N4Igdg9gJgpgziAXAbVABwnAlgFyxMJZABgBpiBdUkANwEMAbAVxiRADUQBfU9TXfIRQBGclVqMWbAOrdeIDNjwEiZYePrNWiDgD0AVHL5LBRUeuqapO6Qe7iYUAObwioAGYAnCAFskAZmocCCQAJktJbRAAHWj3dzseD28-RDIQYKRRCS02WKcnIxAvXyR0zMRwkAY6ACMYBgAFfmUhEE8sJwALHBAI3J1Y+IB9TiTilKygkMRAnOsY6ILR+y4gA
	\begin{center}
		\begin{tikzcd}
			V \arrow[r, "\gg"] \arrow[d, "\ff_V"'] & W \arrow[d, "\ff_W"]   \\
			V^*                                    & W^* \arrow[l, "\ff^*"]
		\end{tikzcd}
	\end{center}
\end{defn}

We can easily generalize this definition so that it applies to more linear isomorphisms than just those between $V$ and $V^*$.

\begin{defn}
	(Natural linear isomorphism).
	
	Let $V$ be a vector space, and let $\FF(V)$ be a vector space constructed from $V$ in such a way that for every vector space $W$ and every linear function $\ff:V \rightarrow W$, there is a linear function $\FF(\ff):\FF(V) \rightarrow \FF(W)$. We say that a linear isomorphism $\ff:V \rightarrow \FF(V)$ is \textit{natural} if and only if the following diagram commutes for all vector spaces $W$ and all linear isomorphisms $\gg:V \rightarrow W$:
	
	% https://tikzcd.yichuanshen.de/#N4Igdg9gJgpgziAXAbVABwnAlgFyxMJZABgBpiBdUkANwEMAbAVxiRADUQBfU9TXfIRQBGclVqMWbAOrdeIDNjwEiZYePrNWiDgD0AVHL5LBRUeuqapO6Qe7iYUAObwioAGYAnCAFskAZmocCCQAJktJbRAAHWj3dzseD28-RDIQYKRRCS02WKcnIxAvXyR0zMRwkAY6ACMYBgAFfmUhEE8sJwALHBAI3J1Y+IB9TiTilKygkMRAnOsY6ILR+y4gA
	\begin{center}
		\begin{tikzcd}
			V \arrow[r, "\gg"] \arrow[d, "\ff_V"'] & W \arrow[d, "\ff_W"]   \\
			\FF(V)                                    & \FF(W) \arrow[l, "\FF(\ff)"]
		\end{tikzcd}
	\end{center}
\end{defn}

\section*{Symmetric and orthogonal linear functions}

The only results from this appendix that are necessary to know are the conditions (3) and (4) of Definition \ref{ch::bilinear_forms_metric_tensors::defn::orthogonal_linear_fn} satisfied by an orthogonal linear function.

\begin{deriv}
\label{ch::appendix::defn::dual_transf_after_id}
    (The adjoint of a linear function with respect to a nondegenerate bilinear form).
    
    Let $V$ and $W$ be finite-dimensional vector spaces, let $B$ be a nondegenerate bilinear form on $V$ and $W$, and consider the musical isomorphisms $\flat_1:V \rightarrow W^{*}$ and $\flat_2:W \rightarrow V^{*}$ induced by $B$.
    
    There is an induced linear map $\ff^\dag:V \rightarrow W$, called the \textit{adjoint of $\ff$ (with respect to $B$)}, that is obtained by using the musical isomorphisms on the domain and codomain of the dual transformation $\ff^*:W^* \rightarrow V^*$. We define $\ff^\dag$ to be the unique map for which the following diagram commutes:
    
    \begin{center}
       % https://tikzcd.yichuanshen.de/#N4Igdg9gJgpgziAXAbVABwnAlgFyxMJZABgBpiBdUkANwEMAbAVxiRADUQBfU9TXfIRRkAjFVqMWbAOoA9AFTdeIDNjwEiI0mOr1mrRBwVK+awZvLi9Uw9O7iYUAObwioAGYAnCAFskZEBwIJC0QBjoAIxgGAAV+dSEQTywnAAscEF1JAxAAHVz3cIyeD28-RFCgpAAmanCo2PjzQ2S0jKz9Nnz3d2MSkC9ff2oqxABmOsjouLMNQwYYd3aJTsN8pycTAbKkCcDgxGquCi4gA
        \begin{tikzcd}
        V \arrow[d, "\flat_1"'] \arrow[r, "\ff^\dag"] & W \arrow[d, "\flat_2"] \\
        W^* \arrow[r, "\ff^*"']                & V^*        
        \end{tikzcd}
    \end{center}
    
    Thus, $\ff^\dag = \flat_2^{-1} \circ \ff^* \circ \flat_1$, or, equivalently, $\flat_2 \circ \ff^\dag = \ff^* \circ \flat_1$.
    
    Since $\flat_2 \circ \ff^\dag = \ff^* \circ \flat_1$, then $\ff^\dag(\vv_0)^{\flat_2} = \ff^*(\vv_0^{\flat_1})$. Using the definition of $\ff^*$ (recall Definition \ref{ch::motivated_intro::defn::dual_transf}), we have $\ff^\dag(\vv_0)^{\flat_2} = \vv_0^{\flat_1} \circ \ff$. This is the same\footnote{Here we've used $B(\cdot, \vv_2)$ to denote the function $\vv_1 \mapsto B(\vv_1, \vv_2)$.} as $B(\cdot, \ff^\dag(\vv_0)) = B(\vv_0, \cdot) \circ \ff$. After evaluating both sides of this most recent equation on $\vv_1 \in V$ and relabeling $\vv_0$ as $\vv_2$, we have:
    
    \begin{align*}
        B(\vv_1, \ff^\dag(\vv_2)) = B(\vv_2, \ff(\vv_1)) \text{ for all $\vv_1, \vv_2 \in V$}.
    \end{align*}
\end{deriv}

Replacing the bilinear form of the above derivation with an inner product, we see that when we have an inner product instead of an arbitrary nondegenerate bilinear form, then $\ff^\dag$ is uniquely characterized by the condition that $\langle \vv_1, \ff^\dag(\vv_2) \rangle = \langle \vv_2, \ff(\vv_1) \rangle$ for all $\vv_1, \vv_2 \in V$. We will use this condition as a means to characterize special classes of linear functions. Specifically, we consider linear functions $\ff$ for which $\ff = \ff^\dag$ and for which $\ff^\dag = \ff^{-1}$. Before investigating such linear functions, we quickly present a definition.

\begin{defn}
    (Transpose of a matrix).
    
    The \textit{transpose} of an $m \times n$ matrix $(a^i{}_j)$ is the $n \times m$ matrix $(a^j{}_i)$. The tranpose of a matrix $\AA$ is denoted $\AA^\top$.
\end{defn}

Now, we consider linear functions $\ff$ for which $\ff = \ff^\dag$ and for which $\ff^\dag = \ff^{-1}$.

\begin{defn} 
\label{ch::bilinear_forms_metric_tensors::defn::symmetric_linear_fn}
    (Symmetric linear function).

    Let $V$ be a vector space, consider a linear function $\ff:V \rightarrow V$. We define $\ff$ to be \textit{symmetric} iff the following equivalent conditions hold:
    
    \begin{enumerate}
        \item $\ff = \ff^\dag$.
        \item $\langle \ff(\vv_1), \vv_2 \rangle = \langle \vv_1, \ff(\vv_2) \rangle$ for all $\vv_1, \vv_2 \in V$.
        \item If $V$ is finite-dimensional and $\hU$ is an orthonormal basis of $V$, then the matrix of $\ff$ relative to $\hU$ is equal to its transpose: $[\ff(\hU)]_{\hU}^\top = [\ff(\hU)]_{\hU}$. (Matrices that are equal to their own transpose are called \textit{symmetric}).
    \end{enumerate}
\end{defn}
        
\begin{proof}
    \mbox{} We need to show that the conditions are indeed equivalent.
    \\ \indent ($1 \iff 2$).
    \\ \indent ($\implies$). Use $\ff = \ff^\dag$ with $\langle \vv_1, \ff^\dag(\vv_2) \rangle = \langle \vv_2, \ff(\vv_1) \rangle$. 
    \\ \indent ($\impliedby$). We have $\langle \ff(\vv_1), \vv_2 \rangle = \langle \vv_1, \ff(\vv_2) \rangle$ for all $\vv_1, \vv_2 \in V$ and $\langle \vv_1, \ff^\dag(\vv_2) \rangle = \langle \vv_2, \ff(\vv_1) \rangle$. Therefore $\langle \vv_1, \ff(\vv_2) \rangle = \langle \vv_1, \ff^\dag(\vv_2) \rangle$. Due to the cancelability of inner products (this follows from the positive-definiteness of inner products), we have $\ff(\vv_2) = \ff^\dag(\vv_2)$ for all $\vv_2 \in V$. So $\ff = \ff^\dag$. 

    ($2 \iff $ 3). 
    \\ \indent ($\implies$). If (2) is true, then the $j$th column of $[\ff(\hU)]_{\hU}$ is $[\ff(\huu_j)]_{\hU}$, so the $^i_j$ entry of $[\ff(\hU)]_{\hU}$ is $\langle \ff(\huu_j), \huu_i \rangle$. Similarly, the $^i_j$ entry of the matrix of $\ff^\dag$ relative to $\hU$ is $\langle 
    \ff^\dag(\huu_j), \huu_i \rangle$. Due to the condition on $\langle \cdot, \cdot \rangle$ induced by the identification $V \cong V^*$, we have $\langle \ff^\dag(\huu_j), \huu_i \rangle = \langle \huu_j, \ff(\huu_i) \rangle = \langle \ff^\dag(\huu_i), \huu_j \rangle$. But this is the $^j_i$ entry of $[\ff(\hU)]_{\hU}$, i.e., the $^i_j$ entry of $[\ff(\hU)]_{\hU}^{\top}$. Thus $[\ff(\hU)]_{\hU} = [\ff(\hU)]_{\hU}^{\top}$.
    \\ \indent ($\impliedby$). Let $\hU = \{\huu_1, ..., \huu_n\}$ be an orthonormal basis for $V$, and let the matrix $\AA$ of $\ff$ relative to $\hU$ satisfy $a_{ij} = a^j_i$. Since $a^i{}_j = \langle \ff(\huu_i), \huu_j \rangle$, then $\langle \ff(\huu_i), \huu_j \rangle = \langle \huu_i, \ff(\huu_j) \rangle$. Extend with multilinearity to obtain the conclusion.
\end{proof}

\begin{defn}
\label{ch::bilinear_forms_metric_tensors::defn::orthogonal_linear_fn}
    (Orthogonal linear function). 
    
    Let $V$ be a vector space over $K$ (where\footnote{This is a very technical condition, and not much attention should be paid to it. We require this so that $2 \neq 0$, which allows us to divide by $2$.} we have $K \neq \Z/2\Z$), consider a linear function $\ff:V \rightarrow V$. We define $\ff$ to be \textit{orthogonal} iff the following equivalent conditions hold:
    
    \begin{enumerate}
        \item $\ff^\dag = \ff^{-1}$.
        \item $\langle \ff(\vv_1), \vv_2 \rangle = \langle \vv_1, \ff^{-1}(\vv_2) \rangle$ for all $\vv_1, \vv_2 \in V$.
        \item $\ff$ preserves inner product: $\langle \vv_1, \vv_2 \rangle = \langle \ff(\vv_1), \ff(\vv_2) \rangle$ for all $\vv_1, \vv_2 \in V$.
        \item $\ff$ preserves length: $\langle \vv, \vv \rangle = \langle \ff(\vv), \ff(\vv) \rangle$ for all $\vv \in V$.
        \item $\ff$ preserves length and angle: $\langle \vv, \vv \rangle = \langle \ff(\vv), \ff(\vv) \rangle$ for all $\vv \in V$ and $\theta(\vv_1, \vv_2) = \theta(\ff(\vv_1), \ff(\vv_2))$ for all $\vv_1, \vv_2 \in V$.
        \item If $V$ is finite-dimensional and $\hU$ is an orthonormal basis of $V$, then the matrix $[\ff(\hU)]_{\hU}$ of $\ff$ relative to $\hU$ has orthonormal columns.
        \item If $V$ is finite-dimensional and $\hU$ is an orthonormal basis of $V$, then $[\ff(\hU)]_{\hU}^{-1} = [\ff(\hU)]_{\hU}^{\top}$.
    \end{enumerate}
\end{defn}
        
\begin{proof}
    We need to show that the conditions are indeed equivalent. To do so, we prove $(3) \iff (4) \iff (5)$ and then $(1) \iff (2) \iff (3) \iff (6) \iff (7)$. The proof of $(3) \iff (4) \iff (5)$ is essentially the same as the proof of Lemma \ref{ch::lin_alg::lemma::orthogonal_linear_fns_preserve_alg_dot_product}; just replace the dot product of that lemma with an inner product. So, it remains to show $(1) \iff (2) \iff (3) \iff (6) \iff (7)$:
    \mbox{} \\
        ($1 \iff 2$).
        \\ \indent ($\implies$). Use $\ff^\dag = \ff^{-1}$ with $\langle \vv_1, \ff^\dag(\vv_2) \rangle = \langle \vv_2, \ff(\vv_1) \rangle$.
        \\ \indent($\impliedby$). $\langle \vv_1, \ff^{-1}(\vv_2) \rangle = \langle \ff(\vv_1), \vv_2 \rangle$ by hypothesis, and $\langle \ff(\vv_1), \vv_2 \rangle = \langle \vv_1, \ff^\dag(\vv_2) \rangle$ by condition on $\langle \cdot, \cdot \rangle$ imposed by identifying $V \cong V^*$ for $\ff^\dag$. Thus $\langle \vv_1, \ff^{-1}(\vv_2) \rangle = \langle \vv_1, \ff^\dag(\vv_2) \rangle$.
        \\ ($2 \iff 3$). Setting $\vv_3 := \ff^{-1}(\vv_2)$ shows that ${(\langle \ff(\vv_1), \vv_2 \rangle = \langle \vv_1, \ff^{-1}(\vv_2) \rangle \text{ for all $\vv_1, \vv_2 \in V$})} \iff$ \\ ${(\langle \ff(\vv_1), \ff(\vv_3) \rangle = \langle \vv_1, \vv_3 \rangle \text{ for all $\vv_1, \vv_3 \in V$})}$.
        \\ ($3 \iff 6$). 
        \\ \indent $(\implies$). We have in particular that $\langle \ff(\huu_i), \ff(\huu_j) \rangle = \langle \huu_i, {\uu}_j \rangle$. Since $\hU$ is orthonormal, $\langle \huu_i, \huu_j \rangle = \delta^i{}_j$. Therefore $\langle \ff(\huu_i), \ff(\huu_j) \rangle = \delta^i{}_j$, so the columns $[\ff(\huu_i)]_{\hU}$ of the matrix of $\ff$ relative to $\hU$ are orthonormal.
        \\ \indent $(\impliedby)$. Since the columns $[\ff(\huu_i)]_{\hU}$ of the matrix of $\ff$ relative to $\hU$ are orthonormal, we have $\langle \ff(\huu_i), \ff(\huu_j) \rangle = \delta^i{}_j$. Extend with multilinearity to obtain the conclusion.
        \\ ($6 \iff 7$).
        \\ \indent $(\implies$). The $^i_j$ entry of $[\ff(\hU)]_{\hU} [\ff(\hU)]_{\hU}^\top$ is $(\text{$i$th row of $[\ff(\hU)]_{\hU}$}) \cdot (\text{$j$th column of $[\ff(\hU)]_{\hU}$})\top = (\text{$i$th row of $[\ff(\hU)]_{\hU}$}) \cdot (\text{$j$th row of $[\ff(\hU)]_{\hU}$}) = \langle \ff(\huu_i), \ff(\huu_j) \rangle = \delta^i{}_j$. Therefore $[\ff(\hU)]_{\hU} [\ff(\hU)]_{\hU}^\top = \II$. A similar argument shows $[\ff(\hU)]_{\hU}^\top [\ff(\hU)]_{\hU} = \II$.
        \\ \indent $(\impliedby)$. Reversing the logic of the forward direction, we know $\langle \ff(\huu_i), \ff(\huu_j) \rangle = \delta^i{}_j$. Therefore (3) is satisfied. Then we use $(3) \implies (6)$.
\end{proof}

\section*{Bubble sort}

\begin{defn}
\label{ch::appendix::defn::bubble_sort}
    (Bubble sort).
    
    Given a permutation $\sigma \in S_n$, the \textit{bubble sort algorithm} operates on the tuple $(\sigma(1), ..., \sigma(n))$, and produces the tuple $(1, ..., n)$ as output.
    
    The following is a description of bubble sort that does not include the popular optimizations:
    
    \begin{itemize}
        \item For each $i \in \{1, ..., n\}$:
        \begin{itemize}
            \item For each $j \in \{1, ..., n - 1\}$, if $\sigma(j) \geq \sigma(j + 1)$, then swap (i.e. apply a transposition to) $\sigma(j)$ and $\sigma(j + 1)$.
        \end{itemize}
    \end{itemize}
    
    This algorithm indeed sorts $(\sigma(1), ..., \sigma(n))$ because the $i$th greatest element of the tuple is in its final position after the $i$th pass of the outer loop. (The first-greatest element of the tuple, or maximum, is in its final position after the first pass of the loop, the second-greatest element of the tuple is in its final position after the second pass of the loop, and so on).
\end{defn}